\section{Results}
\label{sec:results}

% Every number in this section comes from tables/*.tex, which
% scripts/render_paper_tables.py generates from
% results/synthetic_benchmark/summary.json. Do not type numbers by hand.
% LLM and real-data subsections stay as \todo placeholders until the
% pre-registered runs exist.

\subsection{Harness validation on synthetic data}
\label{sec:results-synthetic}

\noindent\fbox{\parbox{\dimexpr\linewidth-2\fboxsep-2\fboxrule}{\small
\textbf{Synthetic data --- not evidence about real markets.} This experiment
validates the search harness on simulated prices with known planted signals
and on null markets with none. Only the two baselines were run; the LLM arm
was not run.}}
\medskip

Table~\ref{tab:synthetic-headline} reports the headline metrics of the
\SynRuns{} planted runs (2 proposers $\times$ 2 SNR levels $\times$
\SynSeeds{} market seeds), Table~\ref{tab:synthetic-recovery} breaks recovery
down by planted signal, Table~\ref{tab:synthetic-null} reports the
\SynNullRuns{} null runs, and Table~\ref{tab:synthetic-oracle} gives the
planted signals' own test-window performance as a ceiling. Every run used
its full budget of \SynBudget{} trials (\SynIncompleteRuns{} incomplete
runs).

\begin{table}[t]
\centering
\caption{Synthetic planted-alpha benchmark, planted markets (mean $\pm$ sd
across complete runs; \SynBudget{} trials per run). Recovery, recall and the
number selected average all runs; the remaining columns describe selected
factors and average only the runs that selected something, so a value
without $\pm$ comes from one run. Composite = equal-weight z-scored
combination of the selected factors on the \SynTestIcDates{}-date sealed
test window. Synthetic data only.}
\label{tab:synthetic-headline}
\resizebox{\textwidth}{!}{% Generated by scripts/render_paper_tables.py from results/synthetic_benchmark/summary.json.
% Do not edit by hand. SYNTHETIC DATA - NOT EVIDENCE ABOUT REAL MARKETS.
\begin{tabular}{llcccccccccc}
\toprule
SNR & Arm & Complete & Recovery & Recall & Selected & True FDP & Test non-sig. & Composite test IC & Composite test ICIR & Behav. novelty & Struct. novelty \\
\midrule
0.05 & evolutionary & 3/3 & 0.08 $\pm$ 0.14 & 0.42 $\pm$ 0.14 & 1.3 $\pm$ 2.3 & 0.00 & 0.50 & 0.0224 & 0.162 & 0.38 & 0.88 \\
0.05 & random & 3/3 & 0.00 $\pm$ 0.00 & 0.42 $\pm$ 0.14 & 0.3 $\pm$ 0.6 & 0.00 & 0.00 & 0.0186 & 0.131 & 0.59 & 0.86 \\
0.15 & evolutionary & 3/3 & 0.33 $\pm$ 0.14 & 0.50 $\pm$ 0.00 & 5.0 $\pm$ 0.0 & 0.00 $\pm$ 0.00 & 0.00 $\pm$ 0.00 & 0.0762 $\pm$ 0.0126 & 0.581 $\pm$ 0.074 & 0.27 $\pm$ 0.07 & 0.85 $\pm$ 0.02 \\
0.15 & random & 3/3 & 0.33 $\pm$ 0.14 & 0.42 $\pm$ 0.14 & 5.0 $\pm$ 0.0 & 0.00 $\pm$ 0.00 & 0.13 $\pm$ 0.12 & 0.0677 $\pm$ 0.0040 & 0.540 $\pm$ 0.050 & 0.35 $\pm$ 0.08 & 0.82 $\pm$ 0.05 \\
\bottomrule
\end{tabular}
}
\end{table}

\begin{table}[t]
\centering
\caption{Recovery of each planted signal: runs in which a \emph{selected}
factor matched it on the test window ($\bar\rho\ge\SynRecoveryThreshold$),
and runs in which \emph{any evaluated trial} matched it on the formation
window ($|\bar\rho|\ge\SynRecoveryThreshold$; ``found'', i.e.\ recall, which
is about as frequent on the null markets). Synthetic data only.}
\label{tab:synthetic-recovery}
\resizebox{\textwidth}{!}{% Generated by scripts/render_paper_tables.py from results/synthetic_benchmark/summary.json.
% Do not edit by hand. SYNTHETIC DATA - NOT EVIDENCE ABOUT REAL MARKETS.
\begin{tabular}{llcccc}
\toprule
SNR & Arm & \texttt{reversal\_5} & \texttt{abnormal\_volume\_20} & \texttt{volume\_return\_corr\_10} & \texttt{drawn\_40} \\
\midrule
0.05 & evolutionary & 0/3 selected, 3/3 found & 1/3 selected, 2/3 found & 0/3 selected, 0/3 found & 0/3 selected, 0/3 found \\
0.05 & random & 0/3 selected, 3/3 found & 0/3 selected, 2/3 found & 0/3 selected, 0/3 found & 0/3 selected, 0/3 found \\
0.15 & evolutionary & 2/3 selected, 3/3 found & 2/3 selected, 3/3 found & 0/3 selected, 0/3 found & 0/3 selected, 0/3 found \\
0.15 & random & 2/3 selected, 3/3 found & 2/3 selected, 2/3 found & 0/3 selected, 0/3 found & 0/3 selected, 0/3 found \\
\bottomrule
\end{tabular}
}
\end{table}

\begin{table}[t]
\centering
\caption{Null markets (SNR $0$, \SynNullSeeds{} seeds): runs that selected
at least one factor, and the selection funnel (mean $\pm$ sd across runs).
Every factor selected here would be a false discovery. Synthetic data only.}
\label{tab:synthetic-null}
\resizebox{\textwidth}{!}{% Generated by scripts/render_paper_tables.py from results/synthetic_benchmark/summary.json.
% Do not edit by hand. SYNTHETIC DATA - NOT EVIDENCE ABOUT REAL MARKETS.
\begin{tabular}{lcccccc}
\toprule
Arm & Complete & Runs selecting $\geq 1$ & Selected & Screen survivors & Confirmed & Composite test $t$ \\
\midrule
evolutionary & 10/10 & 0/10 & 0.0 $\pm$ 0.0 & 4.5 $\pm$ 10.9 & 0.0 $\pm$ 0.0 & n/a \\
random & 10/10 & 0/10 & 0.0 $\pm$ 0.0 & 0.2 $\pm$ 0.6 & 0.0 $\pm$ 0.0 & n/a \\
\bottomrule
\end{tabular}
}
\end{table}

\begin{table}[t]
\centering
\caption{Oracle reference: test-window mean IC (Newey--West $t$) of the planted
signals and of their equal-weight composite, per planted market, and the
share of planted signals that pass BH at \SynTestAlpha{} on the test window
(power). Synthetic data only.}
\label{tab:synthetic-oracle}
\resizebox{\textwidth}{!}{% Generated by scripts/render_paper_tables.py from results/synthetic_benchmark/summary.json.
% Do not edit by hand. SYNTHETIC DATA - NOT EVIDENCE ABOUT REAL MARKETS.
\begin{tabular}{llcccccc}
\toprule
SNR & Seed & \texttt{reversal\_5} & \texttt{abnormal\_volume\_20} & \texttt{volume\_return\_corr\_10} & \texttt{drawn\_40} & Planted composite & Power \\
\midrule
0.05 & 0 & 0.0180 (1.90) & 0.0407 (4.23) & 0.0177 (1.64) & 0.0238 (2.25) & 0.0554 (4.72) & 0.75 \\
0.05 & 1 & 0.0274 (2.59) & 0.0300 (2.56) & 0.0244 (2.28) & 0.0264 (2.54) & 0.0522 (4.65) & 1.00 \\
0.05 & 2 & 0.0160 (1.69) & 0.0424 (3.82) & 0.0258 (2.64) & 0.0308 (3.09) & 0.0554 (4.90) & 1.00 \\
0.15 & 0 & 0.0537 (5.71) & 0.0830 (8.80) & 0.0562 (5.52) & 0.0416 (3.88) & 0.1333 (11.94) & 1.00 \\
0.15 & 1 & 0.0652 (6.39) & 0.0700 (5.86) & 0.0625 (5.92) & 0.0568 (5.12) & 0.1370 (12.16) & 1.00 \\
0.15 & 2 & 0.0506 (5.19) & 0.0870 (7.39) & 0.0701 (7.44) & 0.0681 (7.36) & 0.1445 (12.95) & 1.00 \\
\bottomrule
\end{tabular}
}
\end{table}

We draw six observations from the tables. With \SynSeeds{} seeds per planted
cell and \SynNullSeeds{} null seeds, we read them as descriptions of the
harness and the baselines, not as statistical claims.

\paragraph{The confirmation step, not the formation screen, controls false
selections.} On the null markets neither arm selected a factor
(\SynNullSelecting{evolutionary} runs for GP and \SynNullSelecting{random}
for random search). The formation screen alone was not safe for the
feedback-driven arm: it passed \SynNullMeanSd{evolutionary}{n_screen_survivors}
expressions per null run for GP against
\SynNullMeanSd{random}{n_screen_survivors} for random search, and none of
them was confirmed on validation. GP chooses its next expressions after
seeing formation results, so its formation p-values are biased toward
significance. This is why the protocol treats the formation screen as a
filter and places the claim-carrying multiplicity step on data no proposer
has seen.

\paragraph{The protocol has little power at the low signal-to-noise level.}
At SNR \SynSnrLow{}, GP selected \SynMeanSd{0.05}{evolutionary}{n_selected}
factors per run and random search \SynMeanSd{0.05}{random}{n_selected}; the
share of runs that selected anything was
\SynMean{0.05}{evolutionary}{any_selected} and
\SynMean{0.05}{random}{any_selected}. The confirmation step confirmed
\SynMean{0.05}{evolutionary}{n_confirmed} and
\SynMean{0.05}{random}{n_confirmed} candidates on average, out of
\SynMean{0.05}{evolutionary}{n_screen_survivors} and
\SynMean{0.05}{random}{n_screen_survivors} screen survivors. The planted
composite itself is clearly detectable on the test window at this level
(Table~\ref{tab:synthetic-oracle}), but the individual planted signals are
weaker, and the validation window holds only a fifth of the sessions. The
real-data design therefore uses multi-year windows.

\paragraph{At the higher level, selections are true discoveries but mostly
known ones.} At SNR \SynSnrHigh{} GP selected
\SynMeanSd{0.15}{evolutionary}{n_selected} factors per run and random search
\SynMeanSd{0.15}{random}{n_selected}, and the true FDP was \SynMeanSd{0.15}{evolutionary}{fdp_true} for GP
and \SynMeanSd{0.15}{random}{fdp_true} for random search. The test
non-significance rate, \SynMeanSd{0.15}{evolutionary}{test_nonsignificant_rate}
and \SynMeanSd{0.15}{random}{test_nonsignificant_rate}, therefore measures
power, not error. The composite test IC was
\SynMeanSd{0.15}{evolutionary}{composite_test_ic} (GP) and
\SynMeanSd{0.15}{random}{composite_test_ic} (random), against
\SynOracleHighMin{}--\SynOracleHighMax{} for the planted composite; the
difference between the arms is within one seed standard deviation. Only the
textbook signals were recovered: neither arm selected the library-family
signal (\SynLibraryFamilySelected{} runs) or the drawn signal
(\SynDrawnSelected{} runs), and no evaluated trial matched either
(\SynLibraryFamilyFound{} and \SynDrawnFound{} runs), although both are in
the baselines' search space. At this budget, the part of the benchmark that
prior knowledge cannot solve is untouched by undirected search.

\paragraph{Syntax-based novelty would have misled.} The factors selected at
SNR \SynSnrHigh{} have structural novelty
\SynMeanSd{0.15}{evolutionary}{structural_novelty_mean} (GP) and
\SynMeanSd{0.15}{random}{structural_novelty_mean} (random), which would read
as new structure. Their behavioural novelty is
\SynMeanSd{0.15}{evolutionary}{behavioural_novelty_mean} and
\SynMeanSd{0.15}{random}{behavioural_novelty_mean}: their values are
substantially correlated with those of library factors (behavioural novelty
is one minus the largest such correlation), although only a few are exact
re-spellings. Rediscovery must be measured on values, which is how the LLM
study will report it.

\paragraph{Recall does not measure detection.} A planted expression counts as
``found'' when any evaluated trial matches it on the formation window. On the
null markets, where the same expressions carry no return, recall was
\SynNullMeanSd{evolutionary}{recall_rate} for GP and
\SynNullMeanSd{random}{recall_rate} for random search, against
\SynMeanSd{0.15}{evolutionary}{recall_rate} and
\SynMeanSd{0.15}{random}{recall_rate} on the planted markets at SNR
\SynSnrHigh{}. Recall therefore records that a proposer writes such
expressions (close variants of the short-reversal and abnormal-volume
signals are easy to reach in the grammar), not that it detected a signal; H4 uses recall only as a contrast
with null markets.

\paragraph{Syntactic deduplication undercounts repeated hypotheses.} Genetic
programming produced \SynMeanSd{0.15}{evolutionary}{n_behaviour_duplicates}
evaluated trials per run at SNR \SynSnrHigh{} whose per-date ranks repeated
an earlier trial's, and random search
\SynMeanSd{0.15}{random}{n_behaviour_duplicates}. Counted as separate
hypotheses, such clusters would inflate the formation screen's discoveries,
since a step-up procedure rewards clusters of equal p-values.

\subsection{LLM-guided search on the synthetic benchmark}
\label{sec:results-llm-synthetic}

\todo{Pending: no LLM experiment has been run. After the pre-registered run
(seeds 100--119 $\times$ 2 SNR levels plus null markets, responses recorded
for replay), report the following: the LLM arm in
Tables~\ref{tab:synthetic-headline} to~\ref{tab:synthetic-null}; the H2
McNemar tests on the drawn signal; validity, duplicate and degenerate rates;
recall as a function of trials and of cost, against the compute-matched
arm; served model identifiers; and token usage.}

\subsection{Real-market comparison (RQ1, RQ2)}
\label{sec:results-real}

\todo{Pending: no real-market evaluation has been run. After the
pre-registered reveals, report the following per arm, universe and
replicate: trial accounting and the selection funnel; sealed-test IC, ICIR
and Newey--West $t$ of the composite and of each factor; the test
non-significance rate; long-short net Sharpe with cost sensitivity (0, 5, 10
and 20 bps); turnover; the registry's exploration and reveal counts; and the
H1 tests with their decision.}

\subsection{Memorization and look-ahead (RQ3)}
\label{sec:results-contamination}

\todo{Pending: report the pre/post knowledge-cutoff difference-in-differences
against GP (H3) with the cutoff date's documentation source, the
behavioural-novelty-stratified LLM advantage, and the anonymization-probe
results on the pre-cutoff validation window.}

\subsection{Ablations (RQ4)}
\label{sec:results-ablations}

\todo{Pending: report the feedback on/off test (H4), the rationale on/off and
effort ablations, and the budget-scaling study, on synthetic data and the
validation window only.}
